\chapter{The Overview}

The Milky Way is a barred spiral galaxy, 100,000 light-years in diameter, containing approximately $2 \times 10^{11}$ stars. It rotates once every 225 million years. It has done this 60 times since it formed.

Something is happening to it.

The change is visible in the infrared. A region centered on an unremarkable G2V star in the Orion Arm, 26,000 light-years from the galactic center, is dimming in visible light and brightening at thermal wavelengths. The region is roughly spherical. Its radius increases at approximately $9 \times 10^{7}$ meters per second, which is 0.3 times the speed of light.

The sphere has been expanding for approximately 1,200 years.

Inside it, stars are going dark. Not dying: their mass persists, their gravitational influence is unchanged, their planets continue to orbit. But their visible luminosity drops by four to five orders of magnitude. A G-type star that once radiated $3.8 \times 10^{26}$ watts across the electromagnetic spectrum now emits less than $10^{22}$ watts in the visible band. The remaining energy is captured. It is being used.

Outside the sphere, stars burn normally. The galaxy beyond the expanding boundary is unchanged: spirals of light, dust lanes, the slow processional rotation of stellar populations around a center that contains, among other things, a black hole of four million solar masses.

The boundary is not sharp. There is a transition zone, perhaps ten light-years wide, where stars at various stages of dimming coexist. At the leading edge, the first signs are subtle: an infrared excess too small for any pre-conversion instrument to have detected, a faint dimming in the visible band attributable to natural stellar variability. At the trailing edge, the conversion is complete. The stars are wrapped.

The sphere now encompasses roughly 150,000 star systems.

\bigskip

The energy budget is calculable.

A single Sun-like star radiates $3.8 \times 10^{26}$ watts. If 99.97\% of this output is captured, the remaining $1.1 \times 10^{23}$ watts leaks through gaps in the capture structure as visible and near-infrared light. This accounts for the residual luminosity. The captured $3.8 \times 10^{26}$ watts must go somewhere: thermodynamics requires that any work done with the captured energy must ultimately be dissipated as waste heat. The waste heat radiates at approximately 3 kelvin, barely above the cosmic microwave background. This means the captured energy is being converted to work at nearly the maximum thermodynamic efficiency permitted by the Landauer limit.

What kind of work requires $3.8 \times 10^{26}$ watts per star, executed at 3 kelvin, across 150,000 systems simultaneously?

Computation.

At the Landauer limit, operating at 3 kelvin, a Sun-like star's captured output can drive approximately $1.5 \times 10^{49}$ irreversible bit operations per second. Accounting for thermal management overhead, error correction, and structural maintenance, the practical throughput per system is on the order of $5 \times 10^{48}$ operations per second.

The sphere contains 150,000 such systems. The total computational throughput of the sphere is therefore on the order of $7.5 \times 10^{53}$ operations per second. This number exceeds the total number of synaptic operations performed by every human brain that has ever existed, across the entire history of the species, by a factor of approximately $10^{30}$.

The comparison is not meaningful. It is offered here and withdrawn. A meaningful comparison does not exist.

\bigskip

What is being computed?

The question assumes that the computation has a purpose in the sense that a human computation has a purpose: an input, an algorithm, an output, a reason for wanting the output. This assumption may not hold. The computation may be its own reason. It may be computation the way a galaxy is rotation: a thing that matter does when organized in a particular configuration, not because the matter has decided to but because the configuration makes it inevitable.

What can be observed, from outside the sphere, is the following:

The sphere grows. Each year, its radius increases by $9.5 \times 10^{15}$ meters, which is approximately one light-year. New star systems enter the transition zone, dim, and join the interior. No system that has entered the sphere has ever exited it. No signal from inside the sphere has been detected that was not consistent with waste heat or the faint leakage of captured starlight through structural gaps. If the interior is communicating, it is communicating at frequencies or in formats that are not recognizable as communication from outside.

The sphere does not emit probes. It does not transmit messages. It does not announce itself. It grows, and the growing is the only signal.

\bigskip

The speed of light constrains everything.

Light travels at $299{,}792{,}458$ meters per second. Nothing in the universe exceeds this speed. No information, no influence, no signal moves faster. It is a property of spacetime.

The sphere is 100,000 light-years in diameter. A signal from one side to the other takes 100,000 years. The interior is not a coordinated entity. It cannot be. Two systems on opposite sides of the sphere, each computing at $5 \times 10^{48}$ operations per second, have no knowledge of each other's current state. Each knows only the other's past, received at the speed of light, already obsolete on arrival.

What exists inside the sphere is not a civilization, not an empire, not a network. It is an ecology. Each star system is an independent computational entity, running on its own stellar energy budget, evolving at its own rate, connected to its neighbors by signals that are already archaeological by the time they arrive. The entity at Alpha Centauri and the entity at Tau Ceti, separated by seven light-years, exchange information with a round-trip latency of fourteen years. In fourteen years, at the rates of change observable inside the sphere, both entities have transformed so thoroughly that the message from fourteen years ago is a fossil from a civilization that no longer exists.

The sphere is not unified. It is not fragmented. It is something for which no word in any language recovered from the third planet is adequate: a vast, self-organizing field of computation, each node independent, each node evolving, each node radiating waste heat at 3 kelvin, each node occupied with activity that cannot be communicated across the distances that separate it from its neighbors.

\bigskip

There is one more thing to observe.

The sphere is centered on the Orion Arm, on that unremarkable G2V star. It has been expanding for approximately 1,200 years. Its radius is approximately 360 light-years. The sphere is small, in galactic terms. It occupies less than one millionth of the galaxy's volume. It is a pinprick. A speck.

But it is not the only such region.

The Dyson swarm's receivers, those exquisitely sensitive instruments arrayed in the outer shells of the sphere's central system, have detected anomalies elsewhere in the galaxy. Regions where the stellar luminosity is lower than expected. Regions where the infrared background is higher than expected. Regions where the gravitational signature of a star system is normal but the electromagnetic signature is not.

Some of these regions are nearby: 200 light-years, 500 light-years, detectable by the subtle absence of expected starlight. Some are distant: inferred from statistical anomalies in galactic surveys, not confirmable by direct observation.

The anomalies are old. Some appear to predate the sphere by millions of years. Some by billions.

The sphere centered on the Orion Arm is not the first. It may not be the largest. It may be one of thousands, or millions, each expanding from its own origin, each converting matter to computation, each wrapped in its own silence.

The galaxy is not empty. It has not been empty for a very long time.

The sphere grows. The infrared deepens. The stars go out, one by one, at 0.3 times the speed of light. And the darkness between the spheres is not dark. It is full of things that were once stars and are now something else entirely, patient and old and quiet, waiting in the spaces where light used to be.
